% Generated draft: requires mathematical and source review.
\documentclass[11pt]{article}
\usepackage[margin=25mm]{geometry}
\usepackage{fontspec}
\setmainfont{Latin Modern Roman}
\usepackage{amsmath,amssymb,mathtools,unicode-math}
\setmathfont{Latin Modern Math}
\usepackage{enumitem,needspace,xurl,hyperref}
\setlength{\parindent}{0pt}
\setlength{\parskip}{5pt}
\setlength{\emergencystretch}{3em}
\newcommand{\N}{\mathbb{N}}
\newcommand{\Z}{\mathbb{Z}}
\newcommand{\Q}{\mathbb{Q}}
\newcommand{\R}{\mathbb{R}}
\newcommand{\C}{\mathbb{C}}
\newcommand{\F}{\mathbb{F}}
\newcommand{\A}{\mathbb{A}}
\newcommand{\PP}{\mathbb{P}}
\newcommand{\E}{\mathbb{E}}
\newcommand{\eps}{\varepsilon}
\newcommand{\dd}{\,\mathrm{d}}
\begin{document}
\section*{UnsolvedMath ID: 30003087}
\textbf{Source problem code:} catalogue entry
\subsection{Definitions and mathematical statement}
Let \(d\in\Z\) with \(d\geq2\). A prime power means an integer \(p^r\), where \(p\) is prime and \(r\geq1\) is an integer. Define
\[
q=q(d)=\min\{p^r:d\leq p^{2r}+p^r+1,\ p\text{ prime},\ r\geq1\},\qquad i=q^2+q+1-d.
\]
For this value of \(q\), set
\[
m_1(i)=q+1-i,\qquad m_2(i)=2q+1-i,
\]
\[
\varepsilon_1(i)=\begin{cases}1,&0\leq i\leq q-1,\\0,&\text{otherwise},\end{cases}
\qquad
\varepsilon_2(i)=\begin{cases}1,&i>q+1,\\0,&\text{otherwise},\end{cases}
\]
and
\[
\begin{aligned}
t_{q-1}(i)&=\begin{cases}qi-q^2-q,&i>q+1,\\0,&\text{otherwise},\end{cases}\\
t_q(i)&=\begin{cases}qi,&i\leq q+1,\\2q^2-(i-2)q-1,&i>q+1,\end{cases}\\
t_{q+1}(i)&=\begin{cases}q^2+q-iq,&i\leq q+1,\\0,&i>q+1.\end{cases}
\end{aligned}
\]
Define
\[
h(d)=\frac{\begin{aligned}&q^2+q+1-i-\varepsilon_1(i)m_1(i)-\varepsilon_2(i)m_2(i)\\&\quad-t_{q-1}(i)(q-1)-t_q(i)q-t_{q+1}(i)(q+1)\end{aligned}}{\varepsilon_1(i)+\varepsilon_2(i)+t_{q-1}(i)+t_q(i)+t_{q+1}(i)}.
\]
The supplied conjecture asks whether, for every integer \(d\geq2\), the quantity called the \emph{absolute linear Harbourne constant} \(H(d)\) satisfies
\[
\begin{aligned}
i\leq2q-2&\quad\Longrightarrow\quad H(d)=h(d),\\
i=2q-1&\quad\Longrightarrow\quad H(d)=-\frac{q^3-q^2+2q-2}{q^2+q-1}.
\end{aligned}
\]
No assertion for \(i>2q-1\) is supplied. The record does not define \(H(d)\): in particular, it does not specify the ambient geometry, admissible line configurations, base fields or characteristics, or the extremum and normalization defining this constant. Thus these are the reported conjectural identities, with the meaning of \(H(d)\) requiring clarification from the cited source before the statement is fully self-contained.
\subsection{Short English statement}
The conjecture predicts explicit values of the absolute linear Harbourne constant for the two stated ranges of \(i=q^2+q+1-d\), where \(q\) is the least eligible prime power. The definition and geometric scope of \(H(d)\) are missing from the supplied record. Its partially solved status is supplied dated curation, not an independently verified conclusion.
\subsection{Sources}
\begin{itemize}
\item Local UnsolvedMath registry: supplied complete record for permanent ID \(30003087\), entitled \emph{Explicit Formula for Absolute Linear Harbourne Constants}; exact source problem code \texttt{catalogue entry}. The record belongs to the Oberwolfach Reports Open Problems collection and gives proposed year 2016; no individual proposer is supplied.
\item \emph{Mini-Workshop: Arrangements of Subvarieties, and their Applications in Algebraic Geometry} (2016), Oberwolfach Reports, pp. 653--702, DOI 10.4171/owr/2016/14. Supplied report reference and literature-source label: OWR report 14/2016. \url{https://doi.org/10.4171/owr/2016/14}. The supplied background describes these reports as documenting mathematical workshops and their open-problem discussions.
\item Dumnicki--Harrer--Szpond, supplied literature reference: \url{https://arxiv.org/abs/1507.04080}.
\item Dumnicki--Harrer--Szpond 2018, supplied literature reference: \url{https://doi.org/10.1016/j.ffa.2018.03.001}.
\item Literature assessment checked at \texttt{2026-08-21}: the supplied curation labels the problem partially solved and reports that the displayed formula is verified for all \(d\leq31\) in the stated domain and for \(d=q^2+q+1\) with \(q\) a prime power, while the remaining cases covered by the conjecture are unsettled. This is supplied dated curation derived from the OWR material, not independently verified here; its time-sensitive status and cited claims require independent checking.
\item The registry labels the statement \texttt{corrected\_verified} and reports that the official report passage was checked and rewritten while preserving its mathematical content. This verification description is supplied metadata, not a verification performed for this draft.
\end{itemize}
\end{document}
